\begin{solapartat}
\punts{0,3} Sabem que $\vec{F}_m = -e\vec{v}\times\vec{B}$, llavors com $\vec{B}$ és perpendicular a $\vec{v}$, el mòdul de la força magnètica s'expressa com:
\[ F_m = evB \]

\punts{0,2} Per altre banda la força magnètica és perpendicular a $\vec{v}$, per tant, el resultat serà que aplicarem una força normal constant, és a dir, l'ió descriurà un moviment circular uniforme (mòdul de la velocitat constant).

\punts{0,2} El radi de la trajectòria es pot obtenir a partir de la segona llei de Newton:
\[ \vec{F} = m\vec{a} \]

\punts{0,2} i tenint en compte que l'acceleració centrípeta s'expressa com,
\[ a = \frac{v^2}{r} \]

\punts{0,35} I finalment, $evB = m\,\dfrac{v^2}{r} \Rightarrow r = \dfrac{mv}{eB}$
\end{solapartat}

\begin{solapartat}
\punts{0,2} $m_{12C} = 12m_\mathrm{p} = 2,004\times 10^{-26}\ \mathrm{kg}$.

\punts{0,2} $r_{12C} = \dfrac{0,3}{2} = 0,15\ \mathrm{m}$.

\punts{0,2} $B = \dfrac{mv}{er} = \dfrac{2,004\times 10^{-26}\times 4,80\times 10^5}{1,602\times 10^{-19}\times 0,15} = 0,400\ \mathrm{T}$

\punts{0,2} $m_{14C} = 14m_\mathrm{p} = 2,338\times 10^{-26}\ \mathrm{kg}$.

\punts{0,2} $r_{14C} = \dfrac{m_{14C}v}{eB} = \dfrac{2,338\times 10^{-26}\times 4,80\times 10^5}{1,602\times 10^{-19}\times 0,400} = 0,175\ \mathrm{m}$.

\punts{0,25} els ions $^{14}_{6}\mathrm{C}^-$ surten per una obertura situada a $2\times 0,175 = 0,35\ \mathrm{m}$ de l'entrada.
\end{solapartat}
